\section{Discussion}\label{sec:discussion}
The gain from SG-I is a finite-$(d,n)$ failure probability.  The bulk contribution is not hidden inside an $o_p(1)$ term: its mean is $\tau_0$, its quadratic fluctuation is $q_0$, and its largest-coordinate obstruction is $b_0$.  The cost is equally important.  SG-I excludes many dependent designs, including spatial, image, genomic, and time-series dependence.  Extending the proof requires a quadratic-form concentration property, not merely a sub-Gaussian label.

Relative to representative NRM/CDM results, the present formulation changes four aspects.  It displays explicit remainders, replaces a required power law by trace/Frobenius/operator summaries, states fixed-$n$ targets as random score oracles, and makes sequential multi-spike subtraction visible.  These changes clarify interpretation rather than establishing a universal priority claim.

Important open problems are: a rank-selection oracle inequality for bNRM; dependent Hanson--Wright conditions; repeated-split CDM with dependence between splits; robust sub-Weibull or finite-moment variants; matrix-level bounds sharper in $n$ than the entrywise argument, potentially using matrix concentration tools such as those surveyed by \citet{Tropp2015}; and minimax lower bounds formalizing the fixed-score information limit.
